\documentclass[11pt]{article}
\usepackage[margin=1in]{geometry}
\usepackage{amsmath,amssymb,amsthm}
\usepackage{graphicx}
\usepackage[T1]{fontenc}
\usepackage[utf8]{inputenc}
\usepackage{lmodern}
\usepackage[hidelinks]{hyperref}

\newtheorem{theorem}{Theorem}
\newtheorem{lemma}{Lemma}
\newtheorem{corollary}[theorem]{Corollary}
\newcommand{\C}{\mathbb C}
\newcommand{\R}{\mathbb R}
\newcommand{\Hone}{H^1(\C^+)}
\newcommand{\BMOA}{\operatorname{BMOA}}
\newcommand{\supp}{\operatorname{supp}}

\title{Finite Blaschke Products Answer the Existence Part of a\\
Toeplitz Problem on $H^1(\mathbb C^+)$}
\author{Run \texttt{fa\_banach\_001}, agent \texttt{agent\_lane\_08}}
\date{9 August 2026}

\begin{document}
\maketitle

\section*{Status and source question}

\textbf{Status: substantial partial solution, likely valid, pending human review.}

Bellavita and Peloso \cite{source} study, for a nonconstant inner function
$\Theta$ on the upper half-plane, the codimension-one subspace
\[
 H^1_\Theta=\left\{f\in\Hone:
       \langle f,\Theta\rangle:=\int_{\R}f(x)\overline{\Theta(x)}\,dx=0
                 \right\}
\]
and the restricted Toeplitz operator
$T_{\overline\Theta}:H^1_\Theta\to\Hone$.  Their Question~1 on PDF
page~12 asks whether any inner functions other than $e^{i\tau z}$ yield a
bounded restricted operator and, more broadly, which inner functions do.
Their Theorem~1.3 shows that boundedness is equivalent to
\begin{equation}\label{eq:source-equivalence}
 H^1_\Theta=K^1_\Theta\oplus\Theta H^1,
 \qquad
 K^1_\Theta=\{f\in\Hone:\overline\Theta f\in\overline{\Hone}\}.
\end{equation}

\begin{figure}[ht]
  \centering
  \includegraphics[width=\textwidth]{figures/open_problem_crop.png}
  \caption{Question 1 on source PDF page 12.}
\end{figure}

\section*{Proof intuition}

For a finite Blaschke product $B$ of degree $N$, the quotient
$H^1/BH^1$ is just the $N$-dimensional space of jets at the zeros of $B$.
The single annihilation condition defining $H^1_B$ cuts this quotient down to
dimension $N-1$.

At the endpoint $p=1$, one must take care not to import an unavailable Riesz
projection.  Instead, the model space can be computed directly.  If $Q$ is
the polynomial whose roots are the reflected zeros of $B$, then analytic and
anti-analytic Fourier support force every $f\in K^1_B$ to have the form
$P/Q$.  Boundary integrability is exactly the condition
$\deg P\leq N-2$.  Thus $K^1_B$ also has dimension $N-1$.  The source's
elementary inclusion $K^1_B\oplus BH^1\subseteq H^1_B$ must therefore be an
equality.

\section*{The finite-Blaschke theorem}

\begin{theorem}\label{thm:main}
Let
\[
 B(z)=\eta\prod_{j=1}^{s}
 \left(\frac{z-a_j}{z-\overline{a_j}}\right)^{m_j},
 \qquad a_j\in\C^+,\quad |\eta|=1,
\]
be a nonconstant finite Blaschke product, with distinct $a_j$, positive
multiplicities $m_j$, and degree $N=\sum_jm_j$.  Then
\[
 H^1_B=K^1_B\oplus BH^1.
\]
Consequently $T_{\overline B}:H^1_B\to\Hone$ is bounded.  If
$f=k+Bh$ is the displayed decomposition, its bounded extension satisfies
$T_{\overline B}f=h$.
\end{theorem}

\begin{proof}
We divide the proof into three short steps.

\emph{Step 1: the jet quotient.}
Consider the bounded jet map
\[
 J:\Hone\longrightarrow\C^N,
 \qquad
 Jf=\bigl(f^{(r)}(a_j)\bigr)_{1\leq j\leq s,\ 0\leq r<m_j}.
\]
It is onto.  Indeed, fix $M>N+1$.  Hermite interpolation applied to
functions of the form $P(z)/(z+i)^M$, with $\deg P<N$, realizes arbitrary
prescribed jets; these rational functions belong to $\Hone$.

The kernel of $J$ is $BH^1$.  One inclusion is immediate.  Conversely, if all
the prescribed jets of $f$ vanish, then $f/B$ is analytic in $\C^+$ and has
the same boundary modulus as $f$.  The standard inner--outer factorization
(equivalently, the Smirnov theorem) gives $f/B\in\Hone$.  Hence
\begin{equation}\label{eq:jet-quotient}
 \Hone/BH^1\cong\C^N.
\end{equation}

Let $L_B(f)=\langle f,B\rangle$.  Since every $h\in\Hone$ has integral zero,
\[
 L_B(Bh)=\int_{\R}h(x)\,dx=0,
\]
so $BH^1\subseteq H^1_B=\ker L_B$.  The functional $L_B$ is nonzero: under
the $H^1$--$\BMOA$ duality, it is represented by the nonconstant function
$B$, which is nonzero modulo constants.  It follows from
\eqref{eq:jet-quotient} that
\begin{equation}\label{eq:hyperplane-dimension}
 \dim(H^1_B/BH^1)=N-1.
\end{equation}

\emph{Step 2: the endpoint model space.}
Put
\[
 A(z)=\prod_{j=1}^{s}(z-a_j)^{m_j},
 \qquad
 Q(z)=\prod_{j=1}^{s}(z-\overline{a_j})^{m_j},
 \qquad B=\eta A/Q.
\]
We claim that
\begin{equation}\label{eq:model-description}
 K^1_B=\left\{\frac{P}{Q}: P\text{ is a polynomial and }
                              \deg P\leq N-2\right\}.
\end{equation}

Let $f\in K^1_B$.  There is a $g\in\Hone$ such that
$\overline Bf=\overline g$ on $\R$, and hence
\[
 \overline\eta\,Qf=A\overline g
\]
as tempered distributions on the boundary.  The Fourier transform of $Qf$
is supported in $[0,\infty)$, while that of $A\overline g$ is supported in
$(-\infty,0]$.  Their common Fourier transform is therefore supported at
$\{0\}$.  A tempered distribution with point support is a finite linear
combination of derivatives of the Dirac mass, so $Qf=P$ is a polynomial.
Since $f=P/Q$ belongs to $L^1(\R)$ and $\deg Q=N$, necessarily
$\deg P\leq N-2$.

Conversely, take such a polynomial $P$ and put
$P^\#(z)=\overline{P(\overline z)}$.  Both
\[
 f=P/Q,
 \qquad g=\eta P^\#/Q
\]
belong to $\Hone$, and on $\R$ one has
$\overline Bf=\overline g$.  Thus $f\in K^1_B$, proving
\eqref{eq:model-description}.  In particular,
\begin{equation}\label{eq:model-dimension}
 \dim K^1_B=N-1,
\end{equation}
with $K^1_B=\{0\}$ when $N=1$.

\emph{Step 3: fill the hyperplane.}
If $k\in K^1_B$, write $\overline Bk=\overline g$ with $g\in\Hone$.
Then
\[
 \langle k,B\rangle=\int_{\R}\overline{g(x)}\,dx=0,
\]
so $K^1_B\subseteq H^1_B$.  Also $K^1_B\cap BH^1=\{0\}$: after cancelling
$B$, a boundary function lying in both $H^1$ and $\overline{H^1}$ must be
zero.  Hence $K^1_B$ injects into $H^1_B/BH^1$.  Equations
\eqref{eq:hyperplane-dimension} and \eqref{eq:model-dimension} show that this
injection is onto.  Therefore
$H^1_B=K^1_B\oplus BH^1$.

The boundedness of $T_{\overline B}$ now follows from
\eqref{eq:source-equivalence}.  On $K^1_B$ the Toeplitz operator vanishes,
while $T_{\overline B}(Bh)=h$, which gives the asserted action.
\end{proof}

\begin{corollary}
Every nonconstant rational inner function on $\C^+$ gives a bounded restricted
Toeplitz operator.  In particular,
\[
 B(z)=\frac{z-i}{z+i}
\]
answers the existence part of Question~1 affirmatively; here
$K^1_B=\{0\}$ and $H^1_B=BH^1$.
\end{corollary}

\section*{Verification, scope, and novelty}

The proof treats repeated zeros through the full Hermite jet family.  It also
checks the endpoint feature that a possible $1/x$ rational tail is not in
$L^1$; this is precisely why the numerator degree in
\eqref{eq:model-description} is $N-2$, rather than the familiar $N-1$ from
$K^2_B$.  Since $BH^1$ is closed and $K^1_B$ is finite-dimensional, the
algebraic direct sum is automatically topological, so its coordinate
projections and the resulting Toeplitz extension are bounded.

This is a partial answer to the source's classification problem.  It settles
the existence question and the full rational-inner subcase, but it does not
classify infinite Blaschke products, singular inner functions, or general
meromorphic inner functions.

On 9 August 2026 we searched the run's registry, solution, attempt, and
proof-gap indexes for arXiv:2503.07281 and the core terms.  Bounded web searches
used the exact paper title, the notation $H^1_\Theta$ and $K^1_\Theta$, the
displayed equality, finite Blaschke products, and later work by the source
authors.  They found the source paper and general model-space background, but
no separate work recording this finite-Blaschke answer.  This supports, but
does not guarantee, novelty.

\section*{Human-review focus}

The decisive checks are: (i) the endpoint identity
$\ker J=BH^1$; (ii) the Fourier-support argument proving
\eqref{eq:model-description}; and (iii) the fact that the nonconstant $B$
defines a nonzero $H^1$ functional, producing exactly one dimension of
codimension inside the jet quotient.

\begin{thebibliography}{9}
\bibitem{source}
C. Bellavita and M. M. Peloso,
\emph{On Toeplitz operators on $H^1(\mathbb C^+)$},
arXiv:2503.07281v1 (2025), accepted in Proc. Amer. Math. Soc.

\bibitem{garnett}
J. B. Garnett,
\emph{Bounded Analytic Functions},
Graduate Texts in Mathematics, Springer, 2006.
\end{thebibliography}

\end{document}
